\section{Interpretation and the next useful research steps}

\subsection{What made the saving possible}

The useful structural distinction is between an arithmetic value and a certificate used to establish its sign. A proposed power output does not need to be positive before entering the power argument if its absolute size can be controlled independently. The quotient-radius construction obtains that control from the exact division rows, then spends a fixed margin in the exponent bound to exclude the unwanted Pell branch.

This suggests a systematic procedure for the direct compiler. Build a dependency graph of sign premises; identify terminal proposed values whose sign is used only to obtain a magnitude bound; express those magnitudes polynomially in earlier bounded data; and reserve enough margin below the shared exponent radius. Such a transformation must retain the actual index predicates and prove both branches of the Pell orientation argument. The present nine deletions are a complete verified instance of that procedure, not a claim that every guard can be removed this way.

The coefficient-height construction addresses a different source of apparent complexity. Repeated linear occurrences, rather than individual arithmetic operations, cause many large coefficients in a sum of squares. Private units separate those occurrences while shared product wires keep the residuals quadratic. The complete zero-set bijection makes this a reusable compiler pass, with a transparent variable cost.

\subsection{Tempting shortcuts that remain unproved or fail}

The few-variable branch was examined but did not yield a sound new record in this investigation. In particular, shortening an auxiliary modulus can destroy the strict bound needed to recover a Pell index from a congruence. Removing a radicand factor can destroy the coprimality that turns a rational square product into separate integral squares. Neither proposed numerical saving is included as a result.

Several checks in the direct branch prevent similarly plausible mistakes. The signed-power lemma fails at base one. The quotient radius does not bound all old power outputs linearly. Old and new Pell tuples are not pointwise equal. A three-square guard cannot simply be deleted because the intended semantic value is positive: positivity must be recovered from the retained rows before interpreting that value. The new proofs make these dependencies explicit.

\subsection{What would constitute a further substantial improvement}

A further reduction should establish at least one new complete point in the measured variable/degree/height/support comparison. Three concrete possibilities remain: recover additional quotient signs with a checked index-bootstrap order; reduce the shared arithmetic source before it is compiled; or combine product sharing and coefficient control more efficiently while preserving a fully checked zero-set correspondence. The two retained exact-binomial high quotients need a different early-index argument, so deleting them is not licensed by this report.

The absolute fewest-variable problem is much harder. Classical universal-equation constructions and modern formalizations show that variable count and degree must be reported together.\cite{jones,universalpairs} The present report deliberately does not rename a direct quartic improvement as a solution of that global optimization problem. Its completed output is a family of explicit equations with rigorous local improvements, a complete coefficient-two companion, and a reproducible path to the requested consistency equivalence.
